\chapter{主线之外：独立光子交换的 Poisson--Skellam 随机模型}
\label{app:incoherent-event-chapter}

这一附录不属于正文 PINEM 主方程的“代入特例树”。正文的相干单轨迹响应始终是 $J_n^2(2|\beta|)$；这里主动更换动力学假设，把每一次 $+\hbar\omega$ 与 $-\hbar\omega$ 交换看成独立随机跳跃，因此必须从另一条 birth--death 主方程重新开始。它只能通过弱耦合单光子概率与正文的 $\beta$ 做参数标定，不能被解释成 Bessel 公式的另一种展开。
\label{app:incoherent-event-model}

正文和附录~\ref{app:analytic-series} 始终属于同一条相干 PINEM 分支：单轨迹概率是 \(J_n^2(2|\beta|)\)，随后才对实验标签做统计平均。本附录讨论另一条明确分叉的模型：把每一次 \(+\hbar\omega\) 或 \(-\hbar\omega\) 能量交换视为彼此独立的随机跳跃。这个假设不由随机化 \(\arg\beta\) 得到，也不是一般退相干的普适结论；它对应能量阶梯上的经典 birth--death 过程。

\section{新的主方程：能量阶梯上的对称随机游走}

令 \(p_n(t)\) 表示电子在时刻 \(t\) 位于能量通道
\(
\mathcal E_0+n\hbar\omega
\)
的概率。设吸收一个光子的跳跃率为 \(\Gamma_+\)，发射一个光子的跳跃率为 \(\Gamma_-\)。在短时间 \(\dd t\) 内，只保留一次跳跃，概率守恒给出
\begin{equation}
\boxed{
\frac{\dd p_n}{\dd t}
=\Gamma_+p_{n-1}
+\Gamma_-p_{n+1}
-(\Gamma_++\Gamma_-)p_n.}
\label{eq:birth-death-master}
\end{equation}
这才是本附录真正的新物理输入。它与相干 PINEM 的单位ary相位调制方程不是同一主方程。

若环境或驱动在所考察区间内满足吸收、发射对称
\begin{equation}
\Gamma_+=\Gamma_-=\Gamma,
\label{eq:symmetric-event-rate}
\end{equation}
并取初态
\begin{equation}
p_n(0)=\delta_{n0},
\end{equation}
定义生成函数
\begin{equation}
G(z,t)\equiv\sum_{n=-\infty}^{+\infty}p_n(t)z^n.
\end{equation}
将式~\eqref{eq:birth-death-master} 乘以 \(z^n\) 并对 \(n\) 求和：
\begin{align}
\partial_tG
&=\Gamma\sum_n p_{n-1}z^n
+\Gamma\sum_n p_{n+1}z^n
-2\Gamma\sum_np_nz^n\\
&=\Gamma(z+z^{-1}-2)G.
\end{align}
由 \(G(z,0)=1\)，
\begin{equation}
G(z,t)
=\exp\!\left[\Gamma t(z+z^{-1}-2)\right].
\end{equation}
利用修正 Bessel 函数生成式
\begin{equation}
\exp\!\left[\frac{x}{2}(z+z^{-1})\right]
=\sum_{n=-\infty}^{+\infty}I_n(x)z^n,
\end{equation}
可读出
\begin{equation}
\boxed{
p_n(t)
=\ee^{-2\Gamma t}I_{|n|}(2\Gamma t).}
\label{eq:Skellam-time}
\end{equation}
定义总平均事件数
\begin{equation}
\bar N\equiv2\Gamma t,
\end{equation}
则
\begin{equation}
\boxed{
P_n^{\rm event}(\bar N)
=\ee^{-\bar N}I_{|n|}(\bar N).}
\label{eq:Skellam}
\end{equation}
这就是 Skellam 分布。等价地，也可以把吸收和发射事件数分别看成独立 Poisson 变量
\(
N_\pm\sim\operatorname{Poisson}(\bar N/2)
\)，再取净阶数 \(n=N_+-N_-\)。

\section{怎样与正文的 $\beta$ 接口：只能做弱耦合标定}

正文相干理论在 \(|\beta|\ll1\) 时给出
\begin{equation}
P_{+1}^{\rm coh}=|\beta|^2+O(|\beta|^4),
\qquad
P_{-1}^{\rm coh}=|\beta|^2+O(|\beta|^4).
\end{equation}
事件模型在 \(\bar N\ll1\) 时由式~\eqref{eq:Skellam} 得到
\begin{equation}
P_{\pm1}^{\rm event}
=\frac{\bar N}{2}+O(\bar N^2).
\end{equation}
若希望两个模型共享同一个单光子弱耦合尺度，可以\emph{定义}标定
\begin{equation}
\boxed{\bar N=2|\beta|^2.}
\label{eq:Nbar-incoherent}
\end{equation}
这只是最低阶匹配，不是从相干模型推导出 Skellam 模型。事实上更高阶立即分叉。令 \(x=|\beta|\)，则
\begin{align}
P_1^{\rm coh}
&=J_1^2(2x)=x^2-x^4+O(x^6),\\
P_1^{\rm event}
&=\ee^{-2x^2}I_1(2x^2)=x^2-2x^4+O(x^6),\\
P_2^{\rm coh}
&=\frac{x^4}{4}+O(x^6),\\
P_2^{\rm event}
&=\frac{x^4}{2}+O(x^6).
\end{align}
因此两者只有单光子最低阶相同；多事件统计和干涉结构不同。

还要注意，随机化光学相位并不会把相干概率变成式~\eqref{eq:Skellam}，因为
\begin{equation}
J_n^2(2|\beta|)
\end{equation}
本来就不依赖 \(\arg\beta\)。事件模型要求的是更强的假设：不同单光子交换在概率层而不是振幅层组合。

\section{有限脉冲：把局域事件数代入同一个随机主式}

若仍采用正文短穿越极限下的局域耦合
\begin{equation}
\beta_{\rm loc}(t,\tau)=\beta_0a(t-\tau),
\end{equation}
并沿用弱耦合标定~\eqref{eq:Nbar-incoherent}，则局域总平均事件数取为
\begin{equation}
\boxed{
\bar N(t,\tau)
=2|\beta_0|^2|a(t-\tau)|^2.}
\label{eq:Nbar-local-event}
\end{equation}
若电子到达时间本身是经典统计分布 \(\rho_{\rm e}(t)\)，实验边带概率就是
\begin{equation}
\boxed{
P_n^{\rm event}(\tau)
=\int\dd t\,\rho_{\rm e}(t)
\ee^{-\bar N(t,\tau)}
I_{|n|}\!\left(\bar N(t,\tau)\right).}
\label{eq:event-pulse-average}
\end{equation}
这已经是有限脉冲随机事件模型的完整母式。没有必要为了形式上的“闭式”再同时展开 \(\ee^{-\bar N}\) 和 \(I_n\) 成双重级数；那样只会把一个简单的概率平均重新写成更难读的记账式。

对 Gaussian 包络
\begin{equation}
a(t)=\exp\!\left(-\frac{t^2}{4\sigma_{\rm L}^2}\right),
\end{equation}
局域事件数就是
\begin{equation}
\bar N(t,\tau)
=\bar N_0
\exp\!\left[-\frac{(t-\tau)^2}{2\sigma_{\rm L}^2}\right],
\qquad
\bar N_0=2|\beta_0|^2.
\end{equation}
把这一式直接代入式~\eqref{eq:event-pulse-average} 数值积分即可。

\subsection{弱事件极限}

若 \(\bar N(t,\tau)\ll1\)，则
\begin{equation}
\ee^{-\bar N}I_1(\bar N)
=\frac{\bar N}{2}+O(\bar N^2).
\end{equation}
于是
\begin{equation}
P_{\pm1}^{\rm event}(\tau)
\simeq
|\beta_0|^2
\int\dd t\,\rho_{\rm e}(t)|a(t-\tau)|^2.
\end{equation}
这正好与相干模型的单光子最低阶一致；也再次说明二者的共同部分只到这里为止。

\section{加入横向束斑和孔径}

若局域事件数还依赖横向位置，直接把实验标签扩展为
\(
\lambda=(\bm R_\perp,t)
\)。在经典混合/探测器平均条件下，式~\eqref{eq:event-pulse-average} 推广为
\begin{equation}
\boxed{
P_n^{\rm event,det}
=\frac{\displaystyle
\int_{\mathcal A}\dd^2R_\perp\,W_\perp(\bm R_\perp)
\int\dd t\,W_t(t)
\ee^{-\bar N(\bm R_\perp,t)}
I_{|n|}\!\left(\bar N(\bm R_\perp,t)\right)}
{\displaystyle
\int_{\mathcal A}\dd^2R_\perp\,W_\perp(\bm R_\perp)}.}
\label{eq:event-detector-average}
\end{equation}
若再采用指数横向近场或其他具体场模型，只需把相应的 \(|\beta(\bm R_\perp)|^2\) 代入
\(
\bar N=2|\beta|^2
\)，没有必要另定义一套 \(d_{jk}^n\) 系数。

\section{两条理论分支的关系}

相干 PINEM 的母式是
\begin{equation}
P_n^{\rm coh}=J_n^2(2|\beta|),
\end{equation}
它来自振幅层的相干求和；实验束流的经典平均只是在这之后计算
\(
\int W P_n^{\rm coh}
\)。本附录的事件模型则从式~\eqref{eq:birth-death-master} 开始，单轨迹概率本身已经变成
\begin{equation}
P_n^{\rm event}=\ee^{-\bar N}I_{|n|}(\bar N).
\end{equation}
所以“对不同电子轨迹做经典统计平均”和“光子交换事件本身失去相干”处在两个不同层级。只有先把这两件事分开，才不会把 Skellam 分布误解为正文 Bessel 概率的另一种级数表示。
